\scanpage{302}
\mspage{252}
lemme, \unsure{tout} comme le lemme 1\inserted{(p.~249)}. Je vais considérer $i^*,i_*$ comme des flèches entre les catégories $\mathscr F W_{\mathscr F}^{-1},\mathscr G W_{\mathscr G}^{-1}$. Elles sont induites par les flèches
\diagram{\begin{tikzcd}[ampersand replacement=\&]
\deleted{\unsure{\mathscr G}}\ \mathscr F\arrow[r,shift right,"{i_{0*}}"']\&\mathscr G\arrow[l,shift right,"{i_0^*}"']
\end{tikzcd}}
entre $\mathscr F$ et $\mathscr G$ eux-mêmes
\diagram{\begin{tikzcd}[ampersand replacement=\&,column sep=small,row sep=small]
{i_0^*\colon\mathscr G\mathrel{\unsure{-|\to}}\mathscr F\colon}\&X_1\arrow[d]\&X_3\arrow[l]\arrow[d]\&{\longmapsto}\&X_1\arrow[d]\&\\
\&S\&X_2\arrow[l]\&\&S\&X_2\arrow[l]
\end{tikzcd}}
\diagram{\begin{tikzcd}[ampersand replacement=\&,column sep=small,row sep=small]
{i_{0*}\colon\mathscr F\to\mathscr G\colon}\&X_1\arrow[d]\&\&{\longmapsto}\&X_1\arrow[d]\&X_1\times_S X_2\arrow[l,arrows={-{Bar}}]\arrow[d]\\
\&S\&X_2\arrow[l]\&\&S\&X_2\arrow[l]
\end{tikzcd}.}

Ces deux foncteurs sont adjoints l'un de l'autre, et compatibles aux localiseurs $W_{\mathscr F},W_{\mathscr G}$. Donc \deleted{\unsure{il}}les foncteurs \deleted{\unsure{obtenus}} déduits de $i_0^*,i_{0*}$ par passage aux catégories de fractions sont également adjoints (cf lemme p.~248).

Ainsi, modulo la dém. des lemmes 1 et \unsure{2}, on a prouvé ceci :
% archive sheet number: 302
